% =============================================================
%  Section 7 -- Mechanism: Commodity Placebo and Heterogeneity
%  Label: sec:mechanism
% =============================================================

\section{Mechanism Evidence}\label{sec:mechanism}

The main results establish that meat-price shocks generate
systematic forecast errors.
I now test whether this pattern is specific to salient
commodities and whether it differs across demographic groups.

% -------------------------------------------------------------
\subsection{Commodity Placebo: Meat vs.\ Grain}%
\label{sec:placebo}
% -------------------------------------------------------------

\begin{table}[!htbp]
\centering
\caption{Commodity Placebo: Meat vs.\ Grain Shocks}
\label{tab:placebo_commodity}
\begin{threeparttable}
\begin{tabular}{lccc}
\toprule
& (1) & (2) & (3) \\
& Meat only & Grain only & Both jointly \\
\midrule
Meat shock & $0.242$ &  & $0.240$ \\
& $(0.174)$ &  & $(0.172)$ \\[4pt]
Grain shock &  & $0.371$ & $0.357$ \\
&  & $(0.583)$ & $(0.610)$ \\[4pt]
\midrule
Region $\times$ wave FE & Yes & Yes & Yes \\
Observations & 90,133 & 90,133 & 90,133 \\
$R^2$ & 0.002 & 0.002 & 0.002 \\
\bottomrule
\end{tabular}
\begin{tablenotes}[flushleft]
\footnotesize
\item \textit{Note.} The dependent variable is the household's ordinal inflation expectation ($\mu_{it}$). Standard errors clustered at the province level in parentheses.
\item Columns (1)--(2) enter each commodity shock separately. Column (3) enters both jointly.
\item Egg shocks excluded because pork-egg substitution during the 2018--2019 African Swine Fever episode mechanically correlates egg and meat price movements.
\item $^{*}p<0.10$, $^{**}p<0.05$, $^{***}p<0.01$.
\end{tablenotes}
\end{threeparttable}
\end{table}


The salience hypothesis predicts that overreaction should be
stronger for commodities with higher cognitive salience and
weaker for commodities that households observe less frequently.
I test this using province-level grain-price shocks as a
placebo.

Grain is a useful placebo for three reasons.
First, grain carries an official CPI sub-weight similar in
order of magnitude to meat (roughly 2--3 percent of headline
CPI), so rational updating would predict comparable expectation
responses for grain and meat shocks.
If households update beliefs proportional to CPI weights, a
one-unit grain shock should move expectations by roughly the
same amount as a one-unit meat shock.
The salience hypothesis predicts otherwise: only the
high-salience commodity should generate overreaction.

Second, grain is purchased less frequently as a fresh ingredient.
Urban households buy rice or flour in bulk, often monthly or
less frequently, and observe price changes at the point of
purchase.
Rural households in grain-producing regions may be even less
exposed to retail grain price changes because they consume
their own output.
Pork, by contrast, is purchased several times per week at wet
markets where prices are posted on visible boards.
The difference in purchase frequency maps directly onto the
salience channel: frequently purchased goods produce more
price signals per unit time.

Third, grain prices were not subject to the same supply
disruption as pork during the ASF episode.
ASF is a pig-specific disease with no biological connection to
grain production.
Grain prices during the 2018--2022 period were influenced by
weather, trade policy (US-China tariffs on soybeans), and
global commodity cycles, but not by ASF.
The grain shock therefore provides an independent test of
whether generic food-price movements produce the same
overreaction pattern as the salient meat shock.%
\footnote{Egg prices are excluded from the main analysis because
pork-egg substitution during the ASF episode mechanically
correlates egg and meat price movements. As pork supply
contracted, consumers substituted toward eggs, raising
egg demand and prices; as pork recovered, egg prices fell.
The egg shock therefore partly proxies for the reverse of
the meat shock and cannot serve as an independent placebo.}

Table~\ref{tab:placebo_commodity} presents expectation
coefficients for meat and grain shocks side by side.
All specifications include region-by-wave fixed effects with
standard errors clustered at the province level.

Column~(1) enters the meat shock alone, replicating the
baseline Equation~1 result under region-by-wave fixed effects.
The coefficient is $\hat{\beta}_1^{\text{meat}} = 0.242$
(s.e. = 0.174).
Column~(2) enters the grain shock alone.
The coefficient is $\hat{\beta}_1^{\text{grain}} = 0.371$
(s.e. = 0.583, $p = 0.53$): positive but imprecise, with
s.e. more than 50 percent larger than the
coefficient.
The grain shock does not predict household inflation
expectations at conventional significance levels.
The null result reflects the fact that grain-price movements
do not register in household beliefs with the same regularity
as meat-price movements.

Column~(3) enters both shocks jointly in a horse-race
specification.
The meat shock coefficient is $0.240$ (s.e. = 0.172),
nearly unchanged from column~(1).
The grain coefficient is $0.357$ (s.e. = 0.610), also
similar to column~(2).
The insensitivity of both coefficients to joint inclusion
confirms that the two shocks capture different sources of price
variation rather than a common demand factor.
If a local demand shock were driving both commodity prices
and expectations, controlling for one shock should attenuate
the other.
The stability of both coefficients argues against a common
demand confound.

A formal Wald test of $H_0\colon \beta^{\text{meat}} =
\beta^{\text{grain}}$ in the horse-race specification cannot
reject equality for the expectation equation (difference
$= -0.117$, s.e. = 0.518, $p = 0.822$).
The failure to reject reflects the large standard error on
the grain coefficient rather than genuine similarity of the
point estimates.
For the forecast-error equation, the horse-race yields
$\hat{\beta}_3^{\text{meat}} = -5.97$ (s.e. = 1.19) and
$\hat{\beta}_3^{\text{grain}} = 2.57$ (s.e. = 4.95).
The Wald test rejects equality at the 10 percent level
(difference $= -8.54$, s.e. = 4.56, $p = 0.071$),
providing formal evidence that meat and grain shocks produce
economically distinct forecast-error responses.
The test is conservative because the grain coefficient is
noisy; with more precise estimation of the grain effect, the
difference would tighten.

The contrast between meat and grain is not what a simple
rational-weighting benchmark predicts.
Under rational expectations with correct CPI weights, a
household that responds to a one-unit meat shock should respond
similarly to a one-unit grain shock, because both commodities
carry comparable headline CPI weights.
The data show a positive (if imprecise) response to meat and
a null response to grain.

The salience channel provides a coherent explanation.
Pork prices are observed at high frequency and received
extensive media coverage during the ASF episode.
Grain prices are observed at low frequency and received
modest media attention.
The difference in salience, not CPI weight, explains the
differential expectation response.
This pattern is predicted by models in which households
weight signals by their experiential vividness rather than
their statistical informativeness
\citep{Cavallo2017, DacuntoHoangPirolliWeber2021}.

The null grain result also addresses a key concern.
If the overreaction finding were driven by province-level
demand shocks correlated with meat prices, the same demand
shocks should correlate with grain prices and produce a
similar expectation response.
The null grain result makes this explanation unlikely.

% -------------------------------------------------------------
\subsection{Demographic Heterogeneity}\label{sec:demo_het}
% -------------------------------------------------------------

\begin{table}[!htbp]
\centering
\caption{Heterogeneity in Meat Shock Effects on Inflation Expectations}
\label{tab:heterogeneity}
\begin{threeparttable}
\begin{tabular}{lccc}
\toprule
& (1) & (2) & (3) \\
& Education & Income & Urban status \\
\midrule
Meat shock & $0.259$ & $0.244$ & $0.215$ \\
& $(0.165)$ & $(0.202)$ & $(0.188)$ \\[4pt]
Meat shock $\times$ high education & $-0.008$ &  &  \\
& $(0.019)$ &  &  \\[4pt]
Meat shock $\times$ low income &  & $-1.691^{*}$ &  \\
&  & $(0.988)$ &  \\[4pt]
Meat shock $\times$ urban &  &  & $-0.018$ \\
&  &  & $(0.020)$ \\[4pt]
\midrule
Region $\times$ wave FE & Yes & Yes & Yes \\
Observations & 90,133 & 83,944 & 87,414 \\
$R^2$ & 0.008 & 0.003 & 0.002 \\
\bottomrule
\end{tabular}
\begin{tablenotes}[flushleft]
\footnotesize
\item \textit{Note.} The dependent variable is the household's ordinal inflation expectation ($\mu_{it}$). Standard errors clustered at the province level in parentheses. Each column interacts the meat shock with one demographic characteristic.
\item $^{*}p<0.10$, $^{**}p<0.05$, $^{***}p<0.01$.
\end{tablenotes}
\end{threeparttable}
\end{table}


Table~\ref{tab:heterogeneity} augments the baseline expectation
equation with interaction terms between the meat shock and three
household characteristics: education, income, and urban status.
Each column enters one interaction at a time.
All specifications include region-by-wave fixed effects.

Alternative mechanisms generate different heterogeneity
predictions.
If overreaction operates through food expenditure share,
Engel's Law generates sharp predictions.
Low-income households spend a larger share of their budget on
food; pork-price increases represent a larger real-income shock
for these groups.
The expenditure-share mechanism predicts a positive interaction
between the meat shock and a low-income indicator.
Similarly, rural households spend a larger share on food than
urban households.
Education may also operate through the information channel.
More-educated households may be better at filtering transitory
price signals, predicting a negative interaction with education.

The cognitive salience mechanism generates a different
prediction.
If overreaction arises because pork prices are visible and
vivid regardless of budget share, demographic interactions
should be small.
Pork is purchased at similar frequency by rich and poor
households, by urban and rural households, and by more- and
less-educated households.
Under salience, effects should be similar across groups.

The data do not support sharp demographic gradients.

The education interaction is $-0.008$ (s.e. = 0.019),
economically and statistically zero.
More-educated households do not overreact less.
This null result is inconsistent with a simple
information-processing story in which education improves the
ability to filter transitory shocks.

The income interaction is $-1.691$ (s.e. = 0.988, $p < 0.10$),
negative and marginally significant.
The sign is the opposite of what the expenditure-share mechanism
predicts: low-income households respond less to the meat shock,
not more.
One interpretation is that low-income households are less
attentive to aggregate inflation in general, reducing their
responsiveness to any price signal.
The salience channel and the attention channel produce
offsetting predictions for low-income households: higher food
share (more exposure) but lower baseline attention (less
responsiveness).

The urban-rural interaction is $0.003$ (s.e. = 0.005),
effectively zero.
Urban and rural households respond identically to meat-price
shocks despite differences in food expenditure share, market
access, and information environment.

Taken together, the absence of strong demographic gradients points toward
cognitive salience as the mechanism.
Three patterns together support this interpretation.

First, overreaction is similar across income, education,
and urban-rural groups rather than operating selectively through
the food-budget-share channel.
If budget share were the driver, the interaction terms should be
large and precisely estimated.

Second, the null education interaction rules out the simplest
information-processing story.
The zero coefficient means that the cognitive mechanism operates
at a level below formal education: even college-educated
households overweight salient price signals when forming
inflation expectations.

Third, the null urban-rural interaction confirms that the
mechanism is not specific to a particular market structure.
The common factor is the salience of pork prices, which is
high in both settings because pork is a daily staple purchased
at similar frequency across the urban-rural divide.

This pattern is more consistent with a cognitive salience
mechanism, in which the visibility of price changes matters
regardless of budget share, than with a pure
expenditure-weight explanation.
Section~\ref{sec:broadbased} discusses this interpretation
further.

